\section{Introduction}
\label{section_1}

% The emergence of ChatGPT ushered us into the irreversible era of AI. The major large language models (LLMs) services have already amassed hundreds of millions of monthly users. In response to this sharp increase in demand for LLMs, constructing an efficient LLM inference system has become one of the most significant challenges currently facing the computer architecture community.
The emergence of ChatGPT ushered us into the irreversible era of AI. Today, major Large Language Model (LLM) services serve hundreds of millions of monthly active users, driving an unprecedented demand for efficient inference systems. As a result, designing high-performance LLM inference infrastructure has become one of the central challenges in the computer architecture community.

%LLM inference is fundamentally partitioned into two distinct phases:
% LLM inference can be broadly divided into two distinct phases: prefill and decode. The prefill stage involves the initial processing of the input prompt, which requires simultaneously computing a potentially large number of tokens. This characteristic mandates high parallel processing throughput. Conversely, the decode stage is an autoregressive process, performing computation only for the newly generated token. By reusing intermediate key-value (KV) tensors from the previous iteration, this phase typically exhibits a very low arithmetic intensity.

% In the LLM inference on xPU(e.g., GPU, NPU), a primary performance bottleneck is the memory-bound nature of the decode stage. Contemporary research efforts to mitigate this challenge have focused on two primary strategies. The first involves input batching to facilitate batched LLM inference~\cite{oaken, attacc, orca, duplex, splitwise}. By aggregating multiple input sequences from diverse users, the effective arithmetic intensity of the decode stage is substantially increased. This approach maximizes the utilization of available compute units, enabling simultaneous output generation and enhancing overall system throughput. The second approach is to leverage algorithmic optimization, such as quantization or data pruning(collectively referred to herein as lossy compression)\cite{amove, mx+, tender, microscopiq, qserve, duquant}. Although lossy compression has the disadvantage of degrading model performance, a well-designed lossy compression technique can significantly boost performance by reducing both the required computing resources and memory bandwidth while minimizing the loss of accuracy.

LLM inference can be broadly divided into two distinct phases: prefill and decode. The prefill stage processes the input prompt and requires the simultaneous computation of a large number of tokens, thereby demanding high computational performance. In contrast, the decode stage operates in an autoregressive manner, generating one token at a time. By reusing intermediate Key–Value (KV) tensors from previous steps, this phase typically exhibits low arithmetic intensity and limited parallelism.

On modern xPU platforms (e.g., GPUs and NPUs), the decode stage often becomes the primary performance bottleneck due to its memory-bound nature. To address this challenge, prior work has largely focused on two complementary strategies. The first is input batching, which aggregates multiple user requests to increase effective arithmetic intensity during decoding~\cite{oaken, attacc, orca, duplex, splitwise}. By improving hardware utilization, batching enables higher throughput through concurrent token generation. The second approach involves algorithmic optimizations such as quantization and data pruning—collectively referred to as lossy compression~\cite{amove, mx+, tender, microscopiq, qserve, duquant}. While such techniques inherently introduce some degree of accuracy degradation, carefully designed methods can substantially improve performance by reducing both compute requirements and memory bandwidth, while keeping accuracy loss minimal.

% Especially in GPUs and NPUs, the major challenge has been improving the performance of the decode stage, which generally operates as memory-bound. Recent research has pursued two main directions to address this problem. The first direction is to use input batching and design an inference system for batched LLM. When computing multiple input sequences from multiple users at once, the arithmetic intensity of the decode stage is significantly increased. This allows us to maximize the utilization of the compute units to generate multiple outputs simultaneously, thereby improving the overall system throughput. The second approach is to leverage algorithmic optimization, such as quantization or data pruning~. In this paper, we call these methods lossy compression. 

% A difference in arithmetic intensity between the prefill stage and the decode stage still remains, and recent research addresses this by using heterogeneous systems.

% Lossy compression is powerful method for improving memory-bound performance as it reduces the amount of data loaded from memory, and it also reduces the overall computation volume. However, some compression methods introduce inefficient hardware operations during the decompression or multiply-accumulate process, which is why many lossy compression techniques are often designed specifically for GPU~\cite{AWQ, qserve} architectures or are co-designed with dedicated hardware~\cite{amove, mx+, tender, microscopiq}.


However, two emerging trends in the LLM ecosystem introduce new challenges. The first is the increasing demand for small LLM inference on consumer-grade devices. Traditionally, both academia and industry have focused on datacenter-scale inference systems. Recently, however, open source small LLM models~\cite{Llama3-8B, Qwen3-4B-Thinking-2507, MobileLLM-R1-950M, Gemma3-270m} have demonstrated competitive performance, bringing LLM inference on consumer-grade devices into the spotlight.

% The other shift is the emergence of reasoning LLMs~\cite{Qwen3-8B, Deepseek-R1-Distillated-Llama3-8B}. These models are developed through reinforcement learning using high-quality datasets after pretraining, and they can handle much more complex tasks, such as solving math problems or coding. Recently, even very small models equipped with reasoning capabilities have been released~\cite{Qwen3-4B-Thinking-2507, MobileLLM-R1-950M}.
The second trend is the rise of reasoning LLMs~\cite{Qwen3-8B, Deepseek-R1-Distillated-Llama3-8B}. These models, typically enhanced via reinforcement learning on high-quality datasets, exhibit strong capabilities in complex tasks such as mathematical reasoning and code generation. Notably, reasoning capabilities are no longer limited to large models; even relatively small models have recently begun to incorporate such functionality~\cite{Qwen3-4B-Thinking-2507, MobileLLM-R1-950M}.

% The new problem arises when deploying these high-performance reasoning LLMs within an edge environment. Given that reasoning tasks typically result in the generation of substantially long output sequences, the decode stage dominates the end-to-end latency. However, input batching is infeasible in edge deployment, as usage is generally confined to a single or a limited number of simultaneous users. Consequently, the only remaining optimization avenue is the application of lossy compression.
These trends jointly expose a critical challenge in edge deployment. Reasoning workloads often produce significantly longer output sequences, causing the decode stage to dominate end-to-end latency. At the same time, batching—one of the most effective optimizations in datacenter environments—is largely infeasible on edge devices, where workloads typically involve a single or a small number of concurrent users. As a result, optimization opportunities are effectively limited to algorithmic approaches such as lossy compression.

The problem is that lossy compression suffers from severe performance degradation when applied to reasoning LLMs. Considering that even very small LLMs now support reasoning capabilities, this issue acts as a critical bottleneck that undermines the potential of LLMs on consumer-grade hardware.

% To achieve high-accuracy LLM acceleration in consumer-grade devices, we believe that speculative decoding can become a key solution. Speculative decoding is an algorithm that parallelizes the decode stage by predicting the next tokens to be generated by using smaller draft models. This method can significantly improve the performance of the decode stage while maintaining the quality of the output sentence.

% However, despite the emergence of numerous speculative decoding algorithms, they suffer from issues such as excessive training cost, low performance gain in low-batch, additional memory capacity overhead, and diminished performance improvement with long sequences. These limitations hinder the widespread adoption of speculative decoding in edge LLM inference environments. To address these problems, we propose Cassandra, a hardware-software co-designed speculative decoding method. Our contributions are summarized as follows:

To address this challenge, we argue that speculative decoding offers a promising direction for high-accuracy acceleration in edge environments. Speculative decoding improves decode-stage efficiency by leveraging a lightweight draft model to predict multiple candidate tokens, which are then verified in parallel. This approach enables significant speedup while preserving output quality.

Despite its potential, existing speculative decoding methods face several limitations, including high training cost, limited effectiveness in low-batch settings, additional memory overhead, and diminishing returns for long output sequences. These drawbacks hinder their applicability to edge scenarios. To overcome these challenges, we propose Cassandra, a hardware–software co-designed speculative decoding framework tailored for low-batch inference on consumer-grade devices. Our main contributions are as follows:

\begin{enumerate}[label={(\arabic*)}] % enumerate 환경 시작, label 옵션으로 레이블 형식 지정
    \item 
        We introduce Cassandra, an algorithm–hardware co-designed speculative decoding method. Our key insight is that a compact yet high-quality draft model can be derived directly from the target LLM through fine-grained data selection, without additional training. Based on this observation, Cassandra applies unstructured value pruning and mantissa truncation to the original model, extracting the most salient information at the bit level to construct the draft model. Furthermore, we identified that the large size of the exponent becomes a bottleneck, preventing the performance improvement in this scenario. As two distinct solutions, we propose lossy exponent compression based on the MX format and lossless exponent compression based on unary coding. By incorporating the fine-grained draft model and exponent compression, Cassandra achieves a performance improvement of up to 2.41$\times$ compared to BFloat16 baseline, which is also significantly higher compared to existing methods.
    % training-free speculative decoding
    \item
        In Cassandra, weights and KV cache are stored in a specialized format obtained through unstructured pruning, mantissa truncation, and exponent compression. To execute computation on commercial xPU architectures, these representations must be converted back to standard floating-point format. To prevent this conversion from becoming a performance bottleneck, we propose a lightweight Cassandra encoder–decoder. The proposed design supports both MX-based and unary-based exponent compression within a unified architecture, enabling efficient format conversion with minimal power and area overhead. Furthermore, we provide integration guidelines for incorporating the encoder–decoder into conventional GPU and NPU architectures, along with a memory management scheme that mitigates irregular memory access patterns.
    
    \item 
        To the best of our knowledge, Cassandra is the first self-speculative decoding framework that leverages algorithm–hardware co-design specifically for low-batch LLM inference on consumer-grade edge devices.
        
\end{enumerate}